\subsection{The complete three-class group at the first changed level}

\begin{statement}{A.173}{Proved}
{The degree-eighteen class group, its exact descent, and the next norm operator}
Put \(d=649\), \(\alpha^3=3\), \(\zeta=\zeta_3\), and
\[
 K_1=\Q(\alpha),\quad K_2=\Q(t),\quad t^3=d^2\alpha,
 \qquad F=K_1(\zeta),\quad E_2=K_2(\zeta).
\tag{A.173.1}
\]
These are the actual fields of A.165.  Let \(\sigma\) generate
\(\operatorname{Gal}(E_2/F)\) by \(t\mapsto\zeta t\), and let
\(\tau\) fix \(K_2\) and invert \(\zeta\).
For \(q=11,59\), write \(R_{q,j}\) for the unique prime of
\(E_2\) above the prime of \(F\) at which
\(\alpha\) reduces to \(r_q\zeta^j\), where
\(r_{11}=-2\) and \(r_{59}=-1/6\) in the respective residue fields.
All three residue fields are \(\F_{q^2}\).
With brackets denoting three-primary ideal classes, put
\[
 C=[R_{11,0}],\qquad
 \Delta_q=[R_{q,1}]-[R_{q,2}].
\tag{A.173.2}
\]
The complete class groups and their indicated generators are
\[
 \begin{gathered}
 \operatorname{Cl}(E_2)\{3\}
    =\langle C\rangle\oplus\langle\Delta_{11}\rangle
                         \oplus\langle\Delta_{59}\rangle
    \simeq(\Z/3)^3,\\
 \tau C=C,\qquad \tau\Delta_q=-\Delta_q,\qquad
 \operatorname{Cl}(K_2)\{3\}\simeq\Z/3.
 \end{gathered}
\tag{A.173.3}
\]
The last group is generated by its linear prime above eleven;
ideal extension to \(E_2\) sends that generator to \(C\).
For \(S=\{2,3,11,59,\infty\}\), both corresponding
\(S\)-class groups have zero three-primary part.

The following source classification includes all radicals over
\(E_2\).  Every class in
\((E_2^\times/E_2^{\times3})^{\tau=-1}\) whose actual cubic
Kummer extension becomes everywhere unramified over \(E_3\)
already gives an everywhere unramified extension over \(E_2\).
Its descended class character on \(K_2\) is a multiple of the
character in A.170.  Its pullback character at level three kills
\([P_{2,0,3}]\).

Finally, set \(A_3=\operatorname{Cl}(E_3)\{3\}\), let
\(\sigma_3\) generate \(\operatorname{Gal}(E_3/E_2)\), and
write \(T_3=\sigma_3-1\).  The actual ideal extension and norm
maps give
\[
 j\operatorname{Cl}(E_2)\{3\}\subset3A_3,
 \qquad T_3^2=0\quad\hbox{on }A_3/3A_3.
\tag{A.173.4}
\]
Every assertion concerns the displayed ordinary class groups and
radical source.  The proof does not replace the balanced quotient
of A.155 by an ordinary class group.
\end{statement}

\begin{proof}
\emph{The ramified primes and the complete base-unit norm index.}
The equation in A.173.1 makes \(E_2/F\) cyclic of degree three.
At each of the six indicated primes over eleven and fifty-nine,
the valuation of \(d^2\alpha\) is two.  At the unique three-prime
of \(F\) it is also two.  All seven primes are totally ramified.
There is no other finite ramification: away from three a unit
radical defines an unramified Kummer extension.  All infinite
places are complex.  The prime-to-three computation of A.138
gives \(\operatorname{Cl}(F)\{3\}=0\).

We use the actual units \(U=\alpha^2-2\) and \(\epsilon\) of
A.170.8.  They and \(\zeta\) form a basis of
\(\mathcal O_F^\times/(\mathcal O_F^\times)^3\).
Here is the precise completeness argument.  The unit rank of
\(K_1\) is one, and A.165 proves that \(U\) generates its
pro-three unit completion.  For the quadratic extension
\(F/K_1\), norm and extension satisfy \(Nj=2\) and
\(jN=1+\tau\).  Multiplication by two is invertible modulo
three.  Thus the plus eigenspace is exactly the line of \(U\).
The free minus eigenspace has rank one; A.170 proves the
prime-to-three saturation of \(\epsilon\) by its depth-one
local class, independent of the root-of-unity class.  The
three-primary roots of unity of \(F\) are \(\mu_3\): inclusion
of \(\mu_9\) would identify the degree-six field \(F\) with
the abelian field \(\Q(\mu_9)\), whereas its explicit
automorphisms \(\alpha\mapsto\zeta\alpha\) and \(\tau\)
do not commute.  This proves the asserted three-element basis.

For a nonzero residue \(a\in\F_{q^2}\), define
\(\ell_q(a)\in\F_3\) by
\(a^{(q^2-1)/3}=\zeta^{\ell_q(a)}\).
At a tamely totally ramified cubic place a unit is a local norm
exactly when its residue is a cube.  Necessity follows by
reducing the product of its three conjugates; sufficiency follows
from the surjective norm on principal units, since three is
invertible in the residue field.  Evaluating the displayed units
at the six actual reductions of \(\alpha\) gives
\[
\begin{array}{c|rrr}
 (q,j)&U&\epsilon&\zeta\\ \hline
 (11,0)&0&2&1\\
 (11,1)&2&0&1\\
 (11,2)&1&0&1\\
 (59,0)&0&2&2\\
 (59,1)&1&1&2\\
 (59,2)&2&1&2
\end{array}
\tag{A.173.5}
\]
This matrix has rank three over \(\F_3\).
The exact arithmetic uses \(\zeta^2+\zeta+1=0\), the formulas
for \(\epsilon\) in A.170.8, and repeated squaring in the
two-dimensional residue algebra.  Every cube of a base unit
is a global norm, by taking that unit as a scalar in \(E_2\).
The rank-three lower bound and the three-dimensional upper
bound therefore prove the exact index
\[
 [\mathcal O_F^\times:
       \mathcal O_F^\times\cap N_{E_2/F}E_2^\times]=27.
\tag{A.173.6}
\]
This step needs no local-to-global norm implication.

\emph{Derivation of the ambiguous class count.}
We give the cohomological calculation, including its unit factor.
Write \(D_E,P_E,U_E\) for fractional ideal divisors, principal
divisors and units of \(E=E_2\), and put \(G=\langle\sigma\rangle\).
All maps in the following argument are respectively ideal
extension, divisor, field norm, or the maps induced by these.
Hilbert's theorem 90 applied to
\(1\to U_E\to E^\times\to P_E\to0\) gives
\[
 \begin{gathered}
 P_E^G/P_F\simeq H^1(G,U_E),\\
 H^1(G,P_E)\simeq
       (U_F\cap N E^\times)/N U_E.
 \end{gathered}
\tag{A.173.7}
\]
For the second isomorphism, use the kernel of
\(H^2(G,U_E)=U_F/NU_E\to H^2(G,E^\times)=F^\times/NE^\times\).
The divisor module is a sum of permutation modules and has
\(H^1(G,D_E)=0\).  Taking invariants in
\(0\to P_E\to D_E\to\operatorname{Cl}(E)\to0\) consequently gives
\[
 |\operatorname{Cl}(E)^G|
    =|D_E^G/P_E^G|\,|H^1(G,P_E)|,
 \qquad
 |D_E^G/P_E^G|
    =\frac{h_F\,3^7}{|H^1(G,U_E)|}.
\tag{A.173.8}
\]
The second equality follows directly from
\([D_E^G:D_F]=\prod e_v=3^7\), using the image of ideal
extension for \(D_F\), and then quotienting by \(P_F\).

For completeness, the two-periodic cyclic resolution gives the
unit Herbrand quotient
\[
 \frac{|\widehat H^0(G,U_E)|}{|H^1(G,U_E)|}=\frac13.
\tag{A.173.9}
\]
Indeed the equivariant Dirichlet logarithm identifies the
rational representation of the free unit lattice with
\(\Q[G]^3-\Q\): the nine complex places of \(E\) form three
regular orbits above the complex places of \(F\), and the
sum-of-logs relation removes one trivial representation.
The Herbrand quotients of \(\Z[G]\) and \(\Z\) are respectively
one and three, as the same cyclic resolution shows.
Finite modules have quotient one.  Passing between commensurable
lattices therefore preserves that quotient, as does adding the
finite group of roots of unity.  This proves A.173.9, hence
\(|H^1(G,U_E)|=3[U_F:NU_E]\).
Combining A.173.7--A.173.9 cancels \([U_F:NU_E]\) and yields
\[
 |\operatorname{Cl}(E)^G|
  =\frac{h_F\,3^6}{[U_F:U_F\cap NE^\times]},
 \qquad
 |(\operatorname{Cl}(E)\{3\})^\sigma|=27.
\tag{A.173.10}
\]

\emph{Three independent characters on the indicated ambiguous classes.}
The divisor of \(g_{11}=2+\alpha\) on \(E_2\) is
\(3R_{11,0}\).  The elements \(u_q\) of A.170 have divisor
\(3(R_{q,2}-R_{q,1})\) with the indexing in A.173.2.
Thus \(C,\Delta_{11},\Delta_{59}\) are killed by three and
fixed by \(\sigma\).  The involution fixes \(C\) and interchanges
the two primes in each \(\Delta_q\).

Consider the two additional literal radicals
\[
                 z_1=11(1+6\alpha),\qquad z_2=d.
\tag{A.173.11}
\]
Their cubic extensions of \(E_2\) are everywhere unramified.
At eleven and fifty-nine, all valuations become divisible by
three after the totally ramified extension \(E_2/F\); the
unit Kummer criterion gives unramified extensions there.
At other finite primes different from three they are units.
At two both have residue one, and cubing is an automorphism
on principal units, so these extensions split.
Here are actual local cube checks at three.  First,
\(v_3(d-1)=4>2v_3(3)\), so Hensel's lemma at one makes
\(d\) a cube in \(\Q_3\).  Second,
\[
             (2+\alpha)^3-z_1=6\alpha^2-54\alpha.
\tag{A.173.12}
\]
In \(K_{1,3}\), normalize \(v(\alpha)=1\), so \(v(3)=3\).
The two terms on the right have valuations five and ten.
Put \(\delta=z_1/(2+\alpha)^3-1\) and \(c=\delta/3\).
Then \(v(c)=2\), and the residual of
\(X^3-(1+\delta)\) at \(1+c\) is \(3c^2+c^3\), of valuation
six, while its derivative has valuation three.
One Newton correction has valuation three.  Taylor expansion
at that corrected point makes the next residual have valuation
at least nine, strictly greater than twice the derivative
valuation.  Strict Hensel therefore supplies the cube root in
\(K_{1,3}\).  The two radicals consequently split at three
already over \(F\).

Write \(\chi_i:\operatorname{Cl}(E_2)\{3\}\to\F_3\)
for the arithmetic Frobenius character of \(\sqrt[3]{z_i}\),
with \(\sqrt[3]{z_i}\mapsto\zeta\sqrt[3]{z_i}\).
Each character is \(\sigma\)-invariant: its radical belongs
to \(F\), so the compositum with \(E_2/F\) has abelian
Kummer Galois group.  Since \(\tau(z_i)=z_i\) and
\(\tau(\zeta)=\zeta^{-1}\), both characters are anti-invariant
under \(\tau\).  In particular they kill \(C\).

The Frobenius values on the two differences can be computed
without a class-group algorithm.  At a prime over \(q\),
\(\varpi=t^2/q\) is a uniformizer and
\(\varpi^3=d^4\alpha^2/q^3\).  For \(z_2\) the normalized
unit has residue class \(\alpha^{-2}\) modulo cubes.
At \(j=1,2\) one has \(1+6\alpha=1-\zeta^j\).
For \(z_1\) the normalized units modulo cubes are respectively
\((1-\zeta^j)\alpha^{-2}\) at eleven and
\(1-\zeta^j\) at fifty-nine.  The omitted rational factors
are cubes in \(\F_{q^2}^\times\), since \(3\nmid q-1\).
Taking their cubic residue logarithms and subtracting the
\(j=2\) value from the \(j=1\) value gives
\[
 \begin{array}{c|rr}
       &\Delta_{11}&\Delta_{59}\\ \hline
 \chi_1&0&2\\
 \chi_2&2&1
 \end{array}.
\tag{A.173.13}
\]
The unramified radical \(u\) of A.170 supplies a third character
\(\chi_0\).  It is \(\sigma\)-invariant because \(u\in F\),
and \(\tau\)-invariant by its explicit descended cyclic field.
It therefore kills both \(\Delta_q\).  Its value on \(C\)
is one: the value two on the linear eleven-prime of \(K_2\)
in A.170.4 is multiplied by the actual residue degree two
in \(E_2/K_2\).  Thus the complete pairing matrix is
\[
 \bigl(\chi_i(C),\chi_i(\Delta_{11}),\chi_i(\Delta_{59})\bigr)_{i=0,1,2}
   =\begin{pmatrix}1&0&0\\0&0&2\\0&2&1\end{pmatrix}.
\tag{A.173.14}
\]
It is invertible.  The three indicated ambiguous classes
therefore generate a group of order 27, which equals the
whole invariant group by A.173.10.

\emph{The operator that recovers the entire class group.}
Let \(A=\operatorname{Cl}(E_2)\{3\}\), \(T=\sigma-1\), and
\(N_\sigma=1+\sigma+\sigma^2\).  On actual ideal classes,
\(N_\sigma=jN_{E_2/F}\).  Since
\(\operatorname{Cl}(F)\{3\}=0\), this operator is zero.
Consequently
\[
                       T^2=N_\sigma-3\sigma=-3\sigma.
\tag{A.173.15}
\]
The finite three-group \(A\) makes \(T\) nilpotent:
\(T^{2a}=(-3)^a\sigma^a\) vanishes as soon as \(3^a A=0\).
All three characters in A.173.14 kill \(TA\), by their
\(\sigma\)-invariance.  Their perfect pairing on \(A^\sigma\)
therefore proves \(A^\sigma\cap TA=0\).
A nonzero \(TA\) under a nilpotent operator has a last nonzero
image, contained in \(\ker T\cap TA\).  This contradicts
the preceding zero intersection.  Hence \(TA=0\) and
\(A=A^\sigma\), proving the first assertion of A.173.3.

For \(E_2/K_2\), ideal norm and extension satisfy
\(Nj=2\) and \(jN=1+\tau\).  On three-primary groups this
makes extension injective and identifies its image with
\(A^{\tau=1}\); its inverse there is \(N/2\).
The pairing and eigenspaces already computed give exactly
the cyclic group and generator in A.173.3.
Finally the ordinary-to-\(S\)-class quotient kills each of
the displayed generating prime classes.  Taking three-primary
parts, an exact operation on finite abelian groups, proves
both \(S\)-class assertions.

\emph{All anti-invariant radicals descending from \(E_2\).}
Let \([w]\) be a class in the radical source of the statement.
At every prime outside \(S\), its valuation is divisible by
three: \(E_3/E_2\) is unramified there, so a ramified cubic
Kummer extension cannot become unramified after this base
change.  In the exact \(S\)-divisor sequence, the divisor
\(\operatorname{div}_{S^c}(w)/3\) has a class killed by three.
The proved vanishing of \(\operatorname{Cl}_S(E_2)\{3\}\)
makes that divisor principal.  Dividing by the cube of its
generator gives an actual \(S\)-unit representing \([w]\).
Replacing this representative by \((w/\tau w)^2\) preserves
its class modulo cubes and makes its norm to \(K_2\) exactly
one.  We henceforth use this representative.

Its anti-invariant divisor over eleven and fifty-nine is
\(c_{11}(R_{11,1}-R_{11,2})+
  c_{59}(R_{59,1}-R_{59,2})\).
The other anti-invariant \(S\)-divisors are principal:
at two their two independent differences are the divisors
of \(\sigma(\beta_2)/\sigma^2(\beta_2)\) and of \(u_2\)
from A.170, and the unique three-prime has zero coefficient.
The first pair consists of the other two primes in the
three-element orbit of the divisor of \(\beta_2\); the
second pair is the degree-six pair, obtained by the unramified
degree-three extension of the two other two-primes of \(F\).
Independence of \(\Delta_{11},\Delta_{59}\) in A.173.3
therefore forces \(3\mid c_{11},c_{59}\).
Thus the Kummer extension was unramified at the changed
primes already over \(E_2\).  The extension \(E_3/E_2\) is
unramified at two, so the assumed unramified pullback forces
the original extension to be unramified there as well.

At three the norm-one representative is a unit.  Write
\(k=E_{2,3}\), \(l=E_{3,3}=k(\sqrt[3]{\rho_2})\), with
\(v_k(\rho_2)=2\pmod3\), as in A.162.
The cubic extensions of \(k\) that become unramified over
\(l\) lie in the compositum of \(l\) with the unique
unramified cubic extension of \(k\).
To verify this last assertion directly, a nontrivial pullback
is the unique unramified cubic extension of \(l\), which
is obtained by that unramified base extension; the trivial
pullback is contained in \(l\) itself.
Their Kummer classes consequently lie in the two-dimensional
span of \([\rho_2]\) and the unramified line.
The valuation of \(\rho_2\) is nonzero modulo three, whereas
both the unit \(w\) and the unramified line have valuation
zero.  Its coefficient must therefore vanish.  The original
extension over \(k\) is unramified.  This proves the global
unramified assertion over \(E_2\).

For an actual norm-one representative, the map
\(\widetilde\tau(\sqrt[3]{w})=(\sqrt[3]{w})^{-1}\)
extends \(\tau\), has order two, and commutes with the cubic
Kummer automorphism.  Its fixed field is the descended
cyclic cubic extension of \(K_2\).  Its inertia at any prime
is trivial: inertia of order three would remain after the
quadratic base change to \(E_2\), where it has just been
proved trivial.  The real place has no odd-order inertia.
The complete group \(\operatorname{Cl}(K_2)\{3\}=\Z/3\)
now shows that the character is a multiple of A.170's one.
Under extension to level three, its value on a prime is the
old character on the ideal norm of that prime.  The
distinguished two-prime has residue degree one in
\(K_3/K_2\), so that norm is \(P_{2,0,2}=(\beta_2)\).
Its value is therefore zero, with the precise norm map retained.

\emph{The next ideal extension operator.}
The three generators of \(A\) are classes of totally
ramified primes or differences of such primes.  In
\(E_3/E_2\), ideal extension therefore takes each to
three times its unique upper prime class or difference.
This proves \(jA\subset3A_3\).
The literal equality
\(1+\sigma_3+\sigma_3^2=jN_{E_3/E_2}\) then gives
\(T_3^2=jN_{E_3/E_2}-3\sigma_3\), whose reduction on
\(A_3/3A_3\) is zero.  This is A.173.4.

The finite residue matrices and local cube inequalities used
above are checked in
\path{code/two_inert_level2_genus_pairing_certificate.py}.
The retained output records rank three in A.173.5 and rank two
in A.173.13.  The class-group deductions are the explicit
cohomology, norm and nilpotent-operator arguments given here.
\end{proof}
